% Options for packages loaded elsewhere
\PassOptionsToPackage{unicode}{hyperref}
\PassOptionsToPackage{hyphens}{url}
\documentclass[
]{article}
\usepackage{xcolor}
\usepackage{amsmath,amssymb}
\setcounter{secnumdepth}{-\maxdimen} % remove section numbering
\usepackage{iftex}
\ifPDFTeX
  \usepackage[T1]{fontenc}
  \usepackage[utf8]{inputenc}
  \usepackage{textcomp} % provide euro and other symbols
\else % if luatex or xetex
  \usepackage{unicode-math} % this also loads fontspec
  \defaultfontfeatures{Scale=MatchLowercase}
  \defaultfontfeatures[\rmfamily]{Ligatures=TeX,Scale=1}
\fi
\usepackage{lmodern}
\ifPDFTeX\else
  % xetex/luatex font selection
\fi
% Use upquote if available, for straight quotes in verbatim environments
\IfFileExists{upquote.sty}{\usepackage{upquote}}{}
\IfFileExists{microtype.sty}{% use microtype if available
  \usepackage[]{microtype}
  \UseMicrotypeSet[protrusion]{basicmath} % disable protrusion for tt fonts
}{}
\makeatletter
\@ifundefined{KOMAClassName}{% if non-KOMA class
  \IfFileExists{parskip.sty}{%
    \usepackage{parskip}
  }{% else
    \setlength{\parindent}{0pt}
    \setlength{\parskip}{6pt plus 2pt minus 1pt}}
}{% if KOMA class
  \KOMAoptions{parskip=half}}
\makeatother
\setlength{\emergencystretch}{3em} % prevent overfull lines
\providecommand{\tightlist}{%
  \setlength{\itemsep}{0pt}\setlength{\parskip}{0pt}}
\usepackage{bookmark}
\IfFileExists{xurl.sty}{\usepackage{xurl}}{} % add URL line breaks if available
\urlstyle{same}
\hypersetup{
  hidelinks,
  pdfcreator={LaTeX via pandoc}}

\author{}
\date{}

\begin{document}

\section{Final Red Team Audit:
luxfi/precompile}\label{final-red-team-audit-luxfiprecompile}

\textbf{Date}: 2026-04-13 \textbf{Scope}: 33 canonical precompiles
across all 11 layers \textbf{Commits}: c81f2b3 (red-flagged fixes),
a0f33ee (anchor collision fix), c02cbe1 (VRF), be53b08 (coverage),
06e1fc3 (PRIMER.md) \textbf{Methodology}: Manual code review + static
analysis of all precompile \texttt{Run()} paths, \texttt{RequiredGas()},
state access, and cryptographic operations.

\textbf{Summary}: 3 CRITICAL, 3 HIGH, 4 MEDIUM, 3 LOW, 3 INFO

\begin{center}\rule{0.5\linewidth}{0.5pt}\end{center}

\subsection{CRITICAL FINDINGS}\label{critical-findings}

\subsubsection{{[}CRITICAL{]} C-01: ML-KEM Encapsulate uses crypto/rand
-\/- consensus
split}\label{critical-c-01-ml-kem-encapsulate-uses-cryptorand----consensus-split}

\textbf{File}: \texttt{mlkem/contract.go:201} \textbf{Class}: consensus
\textbf{Adversary capability}: Any EOA sending a transaction
\textbf{Exploit scenario}:

\begin{enumerate}
\def\labelenumi{\arabic{enumi}.}
\tightlist
\item
  User calls ML-KEM precompile at \texttt{0x0200...07} with
  \texttt{OpEncapsulate} (0x01), mode 0x01 (ML-KEM-768), and a valid
  public key.
\item
  \texttt{encapsulate()} at line 201 calls \texttt{pk.Encapsulate()}
  with no seed parameter.
\item
  \texttt{luxfi/crypto/mlkem.PublicKey.Encapsulate()} internally uses
  \texttt{crypto/rand.Reader} to generate the encapsulation randomness.
\item
  Each validator generates different randomness, producing different
  (ciphertext, shared\_secret) pairs.
\item
  Validators disagree on the EVM return value. Consensus split.
\end{enumerate}

This was partially fixed for HPKE (commit c81f2b3 added
\texttt{deriveEffectiveSeed} + \texttt{deterministicReader}) but the
same fix was NOT applied to the ML-KEM precompile. The HPKE precompile
correctly requires a caller-supplied 32-byte seed. ML-KEM does not.

\textbf{Impact}: Chain split on any ML-KEM encapsulate call. Every node
produces a different output. Immediate liveness failure. \textbf{PoC}:
\texttt{cast\ call\ 0x0200...07\ 0x01\ 0x01\ \textless{}1184\_bytes\_valid\_pk\textgreater{}}
-\/- each node returns different 1120 bytes. \textbf{Fix}: Add mandatory
32-byte seed parameter to ML-KEM encapsulate input format (identical
pattern to HPKE). Use \texttt{deriveEffectiveSeed(caller,\ seed)} and
feed the result into
\texttt{mlkem.EncapsulateDeterministic(pk,\ derivedSeed)}. If
\texttt{luxfi/crypto/mlkem} does not expose a deterministic variant, add
one. \textbf{Verified by}: Direct code trace. \texttt{pk.Encapsulate()}
has no seed parameter; the underlying Kyber implementation samples from
\texttt{crypto/rand}.

\begin{center}\rule{0.5\linewidth}{0.5pt}\end{center}

\subsubsection{{[}CRITICAL{]} C-02: X-Wing Encapsulate uses crypto/rand
-\/- consensus
split}\label{critical-c-02-x-wing-encapsulate-uses-cryptorand----consensus-split}

\textbf{File}: \texttt{xwing/contract.go:93} \textbf{Class}: consensus
\textbf{Adversary capability}: Any EOA sending a transaction
\textbf{Exploit scenario}:

\begin{enumerate}
\def\labelenumi{\arabic{enumi}.}
\tightlist
\item
  User calls X-Wing precompile at \texttt{0x2221} with
  \texttt{OpEncapsulate} (0x02) and a valid public key.
\item
  Line 93: \texttt{ct,\ ss,\ err\ :=\ scheme.Encapsulate(pk)} calls
  \texttt{circl/kem/xwing.Scheme().Encapsulate()}.
\item
  circl\textquotesingle s \texttt{Encapsulate} internally uses
  \texttt{crypto/rand} for both the X25519 ephemeral key and the
  ML-KEM-768 encapsulation randomness.
\item
  Each validator generates different \texttt{(ct,\ ss)}. Consensus
  split.
\end{enumerate}

Same class as C-01 but on a different precompile. The X-Wing precompile
was not touched in commit c81f2b3.

\textbf{Impact}: Chain split on any X-Wing encapsulate call.
\textbf{PoC}:
\texttt{cast\ call\ 0x2221\ 0x02\ \textless{}xwing\_pk\_bytes\textgreater{}}
-\/- different output per node. \textbf{Fix}: Add mandatory 32-byte
seed. Use
\texttt{scheme.EncapsulateDeterministic(pk,\ deterministicReader(deriveEffectiveSeed(caller,\ seed)))}
or equivalent. circl supports \texttt{EncapsulateDeterministic} with an
\texttt{io.Reader}. \textbf{Verified by}: Direct code trace.
\texttt{scheme.Encapsulate(pk)} passes nil randomness reader, defaulting
to \texttt{crypto/rand}.

\begin{center}\rule{0.5\linewidth}{0.5pt}\end{center}

\subsubsection{{[}CRITICAL{]} C-03: FHE initTFHE generates keys from
crypto/rand -\/- consensus
divergence}\label{critical-c-03-fhe-inittfhe-generates-keys-from-cryptorand----consensus-divergence}

\textbf{File}: \texttt{fhe/fhe\_ops.go:31-53} \textbf{Class}: consensus
\textbf{Adversary capability}: Any EOA sending the first FHE operation
after node restart \textbf{Exploit scenario}:

\begin{enumerate}
\def\labelenumi{\arabic{enumi}.}
\tightlist
\item
  \texttt{initTFHE()} runs via \texttt{sync.Once} on the first FHE call
  after process start.
\item
  Line 42: \texttt{kg.GenKeyPair()} generates \texttt{secretKey} and
  \texttt{publicKey} using \texttt{crypto/rand}.
\item
  Every validator generates different FHE keys at startup.
\item
  All subsequent FHE operations (encrypt, add, mul, etc.) produce
  ciphertexts under different keys.
\item
  FHE ciphertext handles (stored as \texttt{common.Hash} of ciphertext
  bytes) differ across validators.
\item
  Any transaction that calls \texttt{fhe.add(handle1,\ handle2)} returns
  different results on each node.
\end{enumerate}

Additionally, line 628 in \texttt{tfheGetNetworkPublicKey()} has a
fallback that calls \texttt{rand.Read(result)} if
\texttt{publicKey.MarshalBinary()} fails -\/- another consensus
divergence path.

\textbf{Impact}: Complete consensus failure for all FHE operations.
Every FHE ciphertext is node-specific. \textbf{PoC}: Any FHE
\texttt{asEuint64(42)} call returns a different handle on each node
because the encryption key differs. \textbf{Fix}: FHE keys MUST be
derived from a deterministic seed committed on-chain (e.g., genesis
config or a key-ceremony output stored in chain state).
\texttt{initTFHE()} must accept a \texttt{{[}32{]}byte} master seed and
derive all keys from it using HKDF. The \texttt{rand.Read} fallback in
\texttt{tfheGetNetworkPublicKey} must be removed entirely (return an
error instead). \textbf{Verified by}: Direct code trace.
\texttt{fhe.NewKeyGenerator(params)} followed by \texttt{GenKeyPair()}
with no seed.

\begin{center}\rule{0.5\linewidth}{0.5pt}\end{center}

\subsection{HIGH FINDINGS}\label{high-findings}

\subsubsection{{[}HIGH{]} H-01: bridge/gateway.go and dex/perpetuals.go
use time.Now() in state-modifying
paths}\label{high-h-01-bridgegatewaygo-and-dexperpetualsgo-use-timenow-in-state-modifying-paths}

\textbf{File}: \texttt{bridge/gateway.go:164,194,215,247,293,378,453};
\texttt{dex/perpetuals.go:73,82,393};
\texttt{dex/teleport.go:143,161,435,514};
\texttt{attestation/attestation.go:115,201,366,397,417} \textbf{Class}:
consensus \textbf{Adversary capability}: Any EOA interacting with
bridge, perpetuals, or attestation precompiles \textbf{Exploit
scenario}:

\begin{enumerate}
\def\labelenumi{\arabic{enumi}.}
\tightlist
\item
  \texttt{InitiateBridge()} at gateway.go:164 sets
  \texttt{CreatedAt:\ uint64(time.Now().Unix())}.
\item
  Two validators process the same transaction 50ms apart. Validator A
  sees \texttt{time.Now()} = 1744539600, Validator B sees 1744539601.
\item
  The \texttt{CreatedAt} field propagates into the
  \texttt{BridgeRequest} hash (used as request ID in the
  \texttt{Requests} map).
\item
  If the request ID or any downstream comparison depends on this
  timestamp, validators disagree on state.
\end{enumerate}

Similarly, \texttt{teleport.go:143} uses \texttt{time.Now().UnixNano()}
in \texttt{generateTeleportID()}, making teleport IDs non-deterministic
across validators.

The PRIMER.md at line 225 explicitly claims: "timestamps come from block
time, not \texttt{time.Now()}". This claim is false for these
precompiles.

\textbf{Impact}: Potential consensus split on bridge, teleport, and
perpetuals operations. Severity depends on whether these map values flow
into on-chain state (which they do -\/- they\textquotesingle re stored
in in-memory maps that affect execution outcomes). \textbf{Fix}: Replace
all \texttt{time.Now()} calls in precompile execution paths with
\texttt{accessibleState.GetBlockContext().Timestamp()}. The block
timestamp is consensus-agreed. \textbf{Verified by}: grep for
\texttt{time.Now()} in non-test \texttt{.go} files under precompile/.

\begin{center}\rule{0.5\linewidth}{0.5pt}\end{center}

\subsubsection{{[}HIGH{]} H-02: FHE ciphertext store is in-memory, not
persisted to chain
state}\label{high-h-02-fhe-ciphertext-store-is-in-memory-not-persisted-to-chain-state}

\textbf{File}: \texttt{fhe/contract.go:920-953} \textbf{Class}: state
\textbar{} consensus \textbf{Adversary capability}: Any EOA using FHE
precompile \textbf{Exploit scenario}:

\begin{enumerate}
\def\labelenumi{\arabic{enumi}.}
\tightlist
\item
  Contract A calls \texttt{fhe.asEuint64(42)}, receiving handle
  \texttt{0xabcd...}.
\item
  Handle is stored in \texttt{ctStore.data} (an in-memory
  \texttt{map{[}common.Hash{]}{[}{]}byte}).
\item
  Contract A stores \texttt{0xabcd...} in EVM storage.
\item
  Node restarts. \texttt{ctStore} is empty. FHE keys are regenerated
  (see C-03).
\item
  Contract B reads \texttt{0xabcd...} from storage, calls
  \texttt{fhe.add(0xabcd...,\ 0x1234...)}.
\item
  \texttt{getCiphertext(0xabcd...)} returns \texttt{(nil,\ 0,\ false)}.
\item
  \texttt{performFHEOperation} returns \texttt{common.Hash\{\}} (zero
  hash).
\item
  Silently wrong result: add of two encrypted values returns zero.
\end{enumerate}

This is not just a restart issue. Different nodes joining the network at
different times never have the same \texttt{ctStore} contents.

\textbf{Impact}: FHE precompile is fundamentally broken for multi-node
consensus. Ciphertexts exist only in the local process memory of the
node that created them. \textbf{Fix}: Ciphertext store must be backed by
EVM state (\texttt{SetState}/\texttt{GetState} on the FHE contract
address). Each ciphertext handle maps to a storage slot containing the
serialized ciphertext. This is expensive (SSTOREs) but necessary for
correctness. \textbf{Verified by}: \texttt{ctStore} is a package-level
\texttt{var} initialized with \texttt{make(map{[}...{]})}. No
persistence mechanism.

\begin{center}\rule{0.5\linewidth}{0.5pt}\end{center}

\subsubsection{{[}HIGH{]} H-03: Ring size parsed as uint8 -\/- gas
charged for 255 but ring can trigger O(n) EC
mults}\label{high-h-03-ring-size-parsed-as-uint8----gas-charged-for-255-but-ring-can-trigger-on-ec-mults}

\textbf{File}: \texttt{ring/contract.go:114,173} \textbf{Class}: gas
\textbf{Adversary capability}: Any EOA \textbf{Exploit scenario}:

\begin{enumerate}
\def\labelenumi{\arabic{enumi}.}
\tightlist
\item
  \texttt{RequiredGas()} at line 114 reads \texttt{ringSize} as
  \texttt{int(input{[}2{]})}, a single byte. Maximum 255.
\item
  Gas = \texttt{GasVerifyBase\ +\ 255\ *\ GasVerifyPerMember} =
  \texttt{4000\ +\ 255\ *\ 2500} = \texttt{641,500} gas.
\item
  The actual ring verify loop (line 238-263) does 2 EC scalar
  multiplications + 2 EC additions per ring member on secp256k1.
\item
  For ringSize=255, that is 510 scalar mults + 510 additions, plus 255
  hash-to-point (SVDW map, RFC 9380).
\item
  Each SVDW hash-to-point involves a modular square root
  (\textasciitilde200us on CPU). 255 of them = \textasciitilde51ms.
\item
  The gas cost of 641,500 is far too cheap for 51ms of CPU work. A
  standard \texttt{ecrecover} (one EC operation) costs 3,000 gas.
\item
  An attacker can fill blocks with ring-sig verifications at below-cost
  gas, slowing block production.
\end{enumerate}

\textbf{Impact}: Gas griefing. Attacker can significantly slow block
processing at below-market-rate gas cost. \textbf{Fix}: Increase
\texttt{GasVerifyPerMember} for secp256k1 LSAG to at least 6,000
(reflecting 2 scalar mults per member at 3,000 each, which matches
ecrecover pricing). Consider capping ring size at 64 or 128 in the
precompile itself. \textbf{Verified by}: Manual cost analysis. secp256k1
scalar mult benchmark is \textasciitilde3,000 gas equivalent; LSAG does
2 per member but charges only 2,500.

\begin{center}\rule{0.5\linewidth}{0.5pt}\end{center}

\subsection{MEDIUM FINDINGS}\label{medium-findings}

\subsubsection{{[}MEDIUM{]} M-01: StableSwap Newton iteration with
maxIterations=256 -\/- no gas scaling for
iterations}\label{medium-m-01-stableswap-newton-iteration-with-maxiterations256----no-gas-scaling-for-iterations}

\textbf{File}: \texttt{stableswap/contract.go:52,288} \textbf{Class}:
gas \textbf{Adversary capability}: Any EOA \textbf{Exploit scenario}:

\begin{enumerate}
\def\labelenumi{\arabic{enumi}.}
\tightlist
\item
  StableSwap charges flat \texttt{GasBase\ =\ 5000} gas for all
  operations.
\item
  \texttt{computeD()} runs Newton\textquotesingle s method with up to
  256 iterations (line 288: \texttt{for\ range\ maxIterations}).
\item
  Each iteration performs multiple \texttt{big.Int} multiplications,
  divisions, additions on 256-bit numbers across \texttt{n} tokens.
\item
  With adversarial inputs (e.g., \texttt{amp} close to zero, highly
  imbalanced balances, \texttt{n=16} tokens), convergence can take many
  iterations.
\item
  256 iterations * 16 tokens * 5 big.Int operations \textasciitilde=
  20,480 big.Int ops for 5,000 gas.
\item
  For comparison, MODEXP (precompile 0x05) charges gas proportional to
  modulus size and exponent.
\end{enumerate}

Additionally, \texttt{computeY()} (called from \texttt{getDy}) runs its
own Newton loop of up to 256 iterations, so a \texttt{getDy} call can
trigger 512 total iterations for 5,000 gas.

\textbf{Impact}: Gas griefing on StableSwap. Attacker crafts
pathological inputs that maximize iterations while paying minimal gas.
\textbf{Fix}: Gas should scale with \texttt{n} (number of tokens) and
include a per-iteration component, or cap \texttt{n} at a small value
(e.g., 4) and increase base gas to cover worst-case 256 iterations.
\textbf{Verified by}: Code inspection of \texttt{computeD} and
\texttt{computeY} loops; both iterate up to 256 times.

\begin{center}\rule{0.5\linewidth}{0.5pt}\end{center}

\subsubsection{{[}MEDIUM{]} M-02: HPKE RequiredGas returns 0 for unknown
operations -\/- free gas
consumption}\label{medium-m-02-hpke-requiredgas-returns-0-for-unknown-operations----free-gas-consumption}

\textbf{File}: \texttt{hpke/contract.go:169} \textbf{Class}: gas
\textbf{Adversary capability}: Any EOA \textbf{Exploit scenario}:

\begin{enumerate}
\def\labelenumi{\arabic{enumi}.}
\tightlist
\item
  \texttt{RequiredGas()} returns 0 for any operation byte that is not
  \texttt{OpSingleShotSeal} (0x20).
\item
  \texttt{Run()} at line 182 calls \texttt{RequiredGas()}, gets 0,
  deducts 0 gas.
\item
  \texttt{Run()} then hits the \texttt{default} case at line 200,
  returning an error.
\item
  The error path at line 204 returns \texttt{suppliedGas\ -\ gasCost} =
  \texttt{suppliedGas\ -\ 0} = full gas refund.
\item
  The call consumes zero net gas despite the EVM needing to dispatch,
  deserialize, and process the call.
\end{enumerate}

This is a minor gas accounting issue (the error is returned so no state
changes), but it violates the principle that failed precompile calls
should consume at least the base gas.

The same pattern exists in multiple precompiles:
\texttt{ring/contract.go:95,111}, \texttt{poseidon/contract.go:74},
\texttt{curve25519/contract.go:79}, \texttt{pasta/contract.go:77},
\texttt{babyjubjub/contract.go:100}.

\textbf{Impact}: Negligible economic impact (errors are cheap), but
violates gas accounting invariant. \textbf{Fix}: All
\texttt{RequiredGas()} functions should return a minimum base gas (e.g.,
100) for any non-empty input, even for invalid operations.
\textbf{Verified by}: Code inspection of RequiredGas default branches.

\begin{center}\rule{0.5\linewidth}{0.5pt}\end{center}

\subsubsection{{[}MEDIUM{]} M-03: FHE decrypt exposes plaintext in EVM
return data -\/- breaks confidentiality
model}\label{medium-m-03-fhe-decrypt-exposes-plaintext-in-evm-return-data----breaks-confidentiality-model}

\textbf{File}: \texttt{fhe/contract.go:192} (handleDecrypt selector),
\texttt{fhe/fhe\_ops.go:559-571} \textbf{Class}: crypto
\textbf{Adversary capability}: Any contract that holds a ciphertext
handle \textbf{Exploit scenario}:

\begin{enumerate}
\def\labelenumi{\arabic{enumi}.}
\tightlist
\item
  Contract A stores an FHE-encrypted balance as handle \texttt{H}.
\item
  Any contract (including malicious contract B) that knows handle
  \texttt{H} can call \texttt{fhe.decrypt(H)}.
\item
  \texttt{handleDecrypt} calls \texttt{tfheDecrypt(ct,\ ctType)} which
  uses the global \texttt{decryptor} to decrypt.
\item
  The plaintext is returned as a \texttt{*big.Int} in the EVM return
  data.
\item
  Since EVM return data is visible to the caller (and in trace logs),
  the encrypted value is now public.
\end{enumerate}

There is no access control on which addresses can decrypt which
ciphertext handles. The handle hash itself serves as the "capability",
but handles are 32-byte keccak hashes that can be read from storage
slots by anyone.

\textbf{Impact}: Any FHE-encrypted value can be decrypted by any
contract that knows the handle, destroying the confidentiality guarantee
of FHE. \textbf{Fix}: \texttt{handleDecrypt} should be gated: only the
original encrypting caller (or a threshold of designated decryptors)
should be able to decrypt. Alternatively, remove \texttt{decrypt} from
the on-chain precompile entirely -\/- decryption should happen off-chain
via a threshold decryption protocol, not via an on-chain precompile that
reveals plaintext in calldata. \textbf{Verified by}: Code trace through
\texttt{handleDecrypt} -\textgreater{} \texttt{performFHEDecrypt}
-\textgreater{} \texttt{tfheDecrypt}. No caller authorization check.

\begin{center}\rule{0.5\linewidth}{0.5pt}\end{center}

\subsubsection{{[}MEDIUM{]} M-04: Anchor precompile allows any caller to
submit for any appID -\/- no
authorization}\label{medium-m-04-anchor-precompile-allows-any-caller-to-submit-for-any-appid----no-authorization}

\textbf{File}: \texttt{anchor/contract.go:107-152} \textbf{Class}: state
\textbf{Adversary capability}: Any EOA or contract \textbf{Exploit
scenario}:

\begin{enumerate}
\def\labelenumi{\arabic{enumi}.}
\tightlist
\item
  Legitimate application registers appID \texttt{0xDEAD...} and submits
  anchor at height 1.
\item
  Attacker calls \texttt{submit(0xDEAD...,\ 2,\ 0xFFFF...)} with a
  garbage root hash.
\item
  The anchor precompile stores the garbage root at height 2 for appID
  \texttt{0xDEAD...}.
\item
  The legitimate application can no longer submit at height 2 (already
  taken) and must use height 3+.
\item
  Verifiers who query height 2 get a garbage root, believing the
  application anchored it.
\end{enumerate}

The anchor precompile has no access control. Any address can submit
anchors for any appID. The \texttt{appID} is a 32-byte hash with no
on-chain registration or ownership concept.

\textbf{Impact}: Anchor integrity compromised. Attackers can front-run
legitimate anchors and insert fake roots, breaking the monotonic trust
chain for any CRDT application. \textbf{Fix}: Add an \texttt{owner}
mapping: the first \texttt{submit} for a given \texttt{appID} sets
\texttt{msg.sender} as the owner. Subsequent submits require
\texttt{caller\ ==\ owner}. Alternatively, derive \texttt{appID} from
\texttt{caller} address to prevent cross-caller collision.
\textbf{Verified by}: Code inspection. \texttt{submit()} only checks
\texttt{readOnly} and height monotonicity, not caller identity.

\begin{center}\rule{0.5\linewidth}{0.5pt}\end{center}

\subsection{LOW FINDINGS}\label{low-findings}

\subsubsection{{[}LOW{]} L-01: FROST liftXCache is an
unbounded-in-practice cache with 1024 cap but no
eviction}\label{low-l-01-frost-liftxcache-is-an-unbounded-in-practice-cache-with-1024-cap-but-no-eviction}

\textbf{File}: \texttt{frost/contract.go:148-179} \textbf{Class}: state
\textbf{Adversary capability}: Contract calling FROST verify with many
distinct public keys \textbf{Exploit scenario}:

\begin{enumerate}
\def\labelenumi{\arabic{enumi}.}
\tightlist
\item
  Attacker sends 1024 FROST verify calls with distinct public keys.
\item
  Cache fills to capacity (1024 entries).
\item
  All subsequent FROST verifications with new keys skip the cache (line
  175: \texttt{if\ len(liftXCache)\ \textless{}\ 1024}).
\item
  The cache becomes useless for new keys while holding potentially stale
  entries.
\end{enumerate}

This is a performance issue, not a security vulnerability. The cache is
bounded and thread-safe.

\textbf{Impact}: Performance degradation after cache fills. No
correctness impact. \textbf{Fix}: Use an LRU cache instead of a
fixed-size map with no eviction. \textbf{Verified by}: Code inspection.
Cache has insert-only policy.

\begin{center}\rule{0.5\linewidth}{0.5pt}\end{center}

\subsubsection{{[}LOW{]} L-02: Ed25519 verify returns nil on invalid
input instead of
error}\label{low-l-02-ed25519-verify-returns-nil-on-invalid-input-instead-of-error}

\textbf{File}: \texttt{ed25519/contract.go:88-89} \textbf{Class}:
parsing \textbf{Adversary capability}: None (informational)

When input length != 128, \texttt{Run()} returns
\texttt{(nil,\ remainingGas,\ nil)} -\/- no error. This is the EVM
convention for "verification failed" but differs from other precompiles
(e.g., MLDSA returns an error for invalid input length). Inconsistent
error reporting across precompiles could confuse Solidity developers who
expect \texttt{revert} on malformed input.

\textbf{Impact}: Developer confusion. No consensus or security impact.
\textbf{Fix}: Standardize: either all verify precompiles return nil+nil
on bad input (current Ed25519 behavior) or all return nil+error. Choose
one and apply consistently. \textbf{Verified by}: Comparison of error
handling across ed25519, mldsa, frost, slhdsa, cggmp21.

\begin{center}\rule{0.5\linewidth}{0.5pt}\end{center}

\subsubsection{{[}LOW{]} L-03: MLDSA batch verify gas is decoupled from
per-signature message
length}\label{low-l-03-mldsa-batch-verify-gas-is-decoupled-from-per-signature-message-length}

\textbf{File}: \texttt{mldsa/contract.go:140-147} \textbf{Class}: gas
\textbf{Adversary capability}: Contract submitting batch verify with
large messages

\texttt{requiredGasBatch()} charges
\texttt{BatchVerifyBaseGas\ +\ count\ *\ BatchVerifyPerSigGas} (50k + N
* 40k). It does not account for per-message length. A batch of 10
signatures where each message is 1MB would cost the same gas as 10
signatures with 32-byte messages, despite vastly more hashing work.

\textbf{Impact}: Minor gas undercharging for large-message batch
verifications. \textbf{Fix}: Add per-byte gas component to batch verify:
iterate the header to sum message lengths and add
\texttt{totalMsgBytes\ *\ MLDSAVerifyPerByteGas}. \textbf{Verified by}:
Code inspection of \texttt{requiredGasBatch}.

\begin{center}\rule{0.5\linewidth}{0.5pt}\end{center}

\subsection{INFO FINDINGS}\label{info-findings}

\subsubsection{{[}INFO{]} I-01: PRIMER.md claims "no crypto/rand" but
SYMMETRIC\_CRYPTO\_ANALYSIS.md documents the HPKE
bug}\label{info-i-01-primermd-claims-no-cryptorand-but-symmetric_crypto_analysismd-documents-the-hpke-bug}

PRIMER.md line 344: "No crypto/rand...no non-deterministic library
behavior." SYMMETRIC\_CRYPTO\_ANALYSIS.md line 199 and 268: Documents
the exact HPKE \texttt{crypto/rand} consensus bug and its fix.

The fix was applied to HPKE (commit c81f2b3). But the PRIMER claim is
still false for ML-KEM (C-01), X-Wing (C-02), and FHE (C-03). Update
PRIMER.md to accurately reflect the status.

\begin{center}\rule{0.5\linewidth}{0.5pt}\end{center}

\subsubsection{{[}INFO{]} I-02: GPU fast path in HPKE singleShotSealGPU
is a
no-op}\label{info-i-02-gpu-fast-path-in-hpke-singleshotsealgpu-is-a-no-op}

\textbf{File}: \texttt{hpke/contract.go:431-478}

\texttt{singleShotSealGPU()} creates GPU tensors, calls
\texttt{lattice.KyberEncaps()}, syncs, then... falls through to
\texttt{singleShotSealCPU(params)} at line 477. The GPU result is
computed but discarded. The function always returns the CPU result.

This is not a bug (the CPU path is correct and deterministic), but the
GPU path is dead code that wastes GPU cycles when \texttt{accel} is
available.

\textbf{Fix}: Either wire the GPU KEM result into the HPKE key schedule,
or remove the GPU fast path until it\textquotesingle s properly
integrated.

\begin{center}\rule{0.5\linewidth}{0.5pt}\end{center}

\subsubsection{{[}INFO{]} I-03: BLS12-381 subgroup checks are correctly
implemented}\label{info-i-03-bls12-381-subgroup-checks-are-correctly-implemented}

\texttt{bls12381/contract.go:101} calls \texttt{pt.IsInSubGroup()} for
G1 points after \texttt{IsOnCurve()}. Line 149 does the same for G2
points. This is correct per EIP-2537 and prevents rogue-key and
small-subgroup attacks on pairing inputs. Verified safe.

\begin{center}\rule{0.5\linewidth}{0.5pt}\end{center}

\subsection{VERIFICATION OF PRIOR RED FLAGS (commit
c81f2b3)}\label{verification-of-prior-red-flags-commit-c81f2b3}

\subsubsection{Flag 1: GPU
non-determinism}\label{flag-1-gpu-non-determinism}

\textbf{Status}: VERIFIED SAFE (with caveat). GPU and CPU must produce
identical mathematical results. The comment in
\texttt{mldsa/contract.go:222-228} documents this correctly. However,
the statement "deterministic across validators built the same way" is a
consensus assumption that must be enforced at the operator level (all
validators MUST use the same build tags). This is acceptable.

\subsubsection{Flag 2: Garbage keys}\label{flag-2-garbage-keys}

\textbf{Status}: VERIFIED SAFE. All signature precompiles (mldsa,
slhdsa, frost, cggmp21, ed25519) pass public keys through their
respective library\textquotesingle s \texttt{PublicKeyFromBytes()} or
\texttt{UnmarshalBinary()} which validates key structure. Invalid keys
return errors. Garbage keys do not verify; they do not crash.

\subsubsection{Flag 3: HPKE seed
collision}\label{flag-3-hpke-seed-collision}

\textbf{Status}: FIX VERIFIED.
\texttt{deriveEffectiveSeed(caller,\ raw)} at
\texttt{hpke/contract.go:384-392} correctly domain-separates with
\texttt{"HPKE\_SEAL\_v1"\ \textbar{}\textbar{}\ caller.Bytes()\ \textbar{}\textbar{}\ raw{[}:{]}}.
Two different callers with the same raw seed produce different effective
seeds. Same caller with same seed is deterministic (consensus safe).

\textbf{Attempted bypass}: Could an attacker force two contracts to
share a caller address? No. The \texttt{caller} parameter comes from the
EVM execution context (\texttt{msg.sender}), which is the direct caller
of the precompile. Two different contracts always have different
addresses. A delegatecall from proxy A to implementation B would use
A\textquotesingle s address as caller, which is correct (the proxy is
the entity choosing the seed).

\subsubsection{Flag 4: init()
double-register}\label{flag-4-init-double-register}

\textbf{Status}: FIX VERIFIED. \texttt{modules/registerer.go:219-230}
now produces descriptive error messages including both the collision key
and address. The fail-fast behavior was already present (returning
error); the fix improves diagnostics.

\subsubsection{Flag 5: isLegacyFormat
heuristic}\label{flag-5-islegacyformat-heuristic}

\textbf{Status}: FIX VERIFIED. \texttt{mldsa/contract.go:425-429}
confirms both \texttt{isLegacyFormat()} and \texttt{RunLegacy()} are
deleted. The comment explains the removal. All ML-DSA input now requires
an explicit mode byte (0x44, 0x65, or 0x87). No bypass possible -\/- the
dispatch in \texttt{Run()} at line 179 only accepts
\texttt{OpBatchVerify} (0x10) or a valid mode byte.

\begin{center}\rule{0.5\linewidth}{0.5pt}\end{center}

\subsection{VERIFIED SAFE}\label{verified-safe}

The following components were examined and found to be correctly
implemented:

\begin{itemize}
\tightlist
\item
  \textbf{BLS12-381 (bls12381/)}: Subgroup checks on G1 and G2 inputs.
  EIP-2537 compliant. No rogue-key attack surface.
\item
  \textbf{Ed25519 (ed25519/)}: Standard library
  \texttt{crypto/ed25519.Verify()} for CPU path. Fixed-size 128-byte
  input prevents buffer issues.
\item
  \textbf{Poseidon (poseidon/)}: gnark-crypto \texttt{poseidon2.Hash()}
  is deterministic. No randomness. Field element validation via
  \texttt{fr.Element.SetBytes()} which reduces modulo the BN254 scalar
  field order.
\item
  \textbf{Pedersen (pedersen/)}: gnark-crypto point operations are
  deterministic. Commitment binding/hiding properties depend on the DLOG
  assumption.
\item
  \textbf{Blake3 (blake3/)}: zeebo/blake3 is deterministic. Keyed hash
  and Merkle modes correctly implemented.
\item
  \textbf{secp256r1 (secp256r1/)}: EIP-7212 P-256 verify. Standard ECDSA
  verification. No key generation on-chain.
\item
  \textbf{sr25519 (sr25519/)}: Verify-only. ChainSafe go-schnorrkel.
  Deterministic.
\item
  \textbf{VRF (vrf/)}: Verify-only (no proving). RFC 9381
  ECVRF-EDWARDS25519-SHA512-ELL2. Correct domain separation tags.
  Hash-to-curve uses elligator2 (Edwards25519 native, not
  try-and-increment).
\item
  \textbf{KZG4844 (kzg4844/)}: Uses go-kzg-4844 trusted setup.
  Deterministic polynomial evaluation and proof verification.
\item
  \textbf{Anchor (anchor/)}: Deterministic keccak256 storage slot
  derivation. Monotonic height enforcement correct. (Auth issue flagged
  as M-04.)
\item
  \textbf{CGGMP21 (cggmp21/)}: Verify-only precompile. No threshold
  signing on-chain.
\item
  \textbf{Ringtail (ringtail/)}: Verify-only. Lattice signature
  verification is deterministic.
\item
  \textbf{Pasta (pasta/)}: Pallas/Vesta point operations via
  gnark-crypto. Deterministic.
\item
  \textbf{BabyJubJub (babyjubjub/)}: BN254 twisted Edwards operations.
  Deterministic.
\item
  \textbf{Curve25519 (curve25519/)}: Point operations on Edwards25519
  via filippo.io/edwards25519. Deterministic.
\item
  \textbf{X25519 (x25519/)}: Diffie-Hellman verify-only (scalar mult).
  No key generation.
\item
  \textbf{Quasar (quasar/)}: BLS aggregate, Verkle, Ringtail, hybrid
  verification. All deterministic verify-only.
\item
  \textbf{Stableswap Newton convergence}: \texttt{maxIterations=256}
  with \texttt{\textbar{}d\ -\ prevD\textbar{}\ \textless{}=\ 1}
  termination is deterministic (same inputs always take the same number
  of iterations with \texttt{big.Int} arithmetic). Gas undercharging is
  M-01 but no consensus risk.
\end{itemize}

\begin{center}\rule{0.5\linewidth}{0.5pt}\end{center}

\subsection{Blue Handoff}\label{blue-handoff}

\textbf{What Blue got right}:

\begin{itemize}
\tightlist
\item
  HPKE domain separation fix (c81f2b3) is cryptographically sound. The
  \texttt{deriveEffectiveSeed} construction is correct.
\item
  Legacy format removal in MLDSA eliminates an ambiguous parsing path.
\item
  BLS12-381 subgroup checks are complete and correct per EIP-2537.
\item
  Anchor collision fix (a0f33ee) correctly separates FHE and anchor
  address ranges.
\item
  Ring signature hash-to-point uses SVDW (RFC 9380) with proper domain
  separation tag.
\item
  Module registerer error messages are now diagnostic enough to catch
  collisions at init time.
\end{itemize}

\textbf{What Blue missed}:

\begin{itemize}
\tightlist
\item
  ML-KEM and X-Wing have the same \texttt{crypto/rand} consensus bug
  that was fixed in HPKE.
\item
  FHE subsystem is fundamentally non-deterministic (keygen from rand,
  in-memory ciphertext store).
\item
  \texttt{time.Now()} in 4 precompile packages (bridge, dex/perpetuals,
  dex/teleport, attestation).
\item
  No authorization model for anchor or FHE decrypt.
\item
  Ring signature gas undercharging relative to actual EC operation cost.
\end{itemize}

\textbf{Fix priority for Blue}:

\begin{enumerate}
\def\labelenumi{\arabic{enumi}.}
\tightlist
\item
  \textbf{C-01 + C-02}: Add deterministic seed to ML-KEM and X-Wing
  encapsulate (same pattern as HPKE).
\item
  \textbf{C-03 + H-02}: Redesign FHE key management (deterministic
  keygen from chain state) and ciphertext persistence (EVM state-backed
  store).
\item
  \textbf{H-01}: Replace all \texttt{time.Now()} in precompile execution
  paths with block timestamp.
\item
  \textbf{H-03}: Increase ring signature per-member gas to match EC
  operation cost.
\item
  \textbf{M-01}: Add iteration-aware gas to StableSwap.
\item
  \textbf{M-03}: Gate FHE decrypt with caller authorization.
\item
  \textbf{M-04}: Add appID ownership to anchor.
\end{enumerate}

\textbf{Re-review scope}: C-01, C-02, C-03, H-01, H-02 after fixes.

\begin{center}\rule{0.5\linewidth}{0.5pt}\end{center}

RED COMPLETE. Findings ready for Blue. Total: 3 critical, 3 high, 4
medium, 3 low, 3 info Top 3 for Blue to fix:

\begin{enumerate}
\def\labelenumi{\arabic{enumi}.}
\tightlist
\item
  ML-KEM + X-Wing encapsulate use crypto/rand (consensus split)
\item
  FHE keygen from crypto/rand + in-memory ciphertext store (consensus
  split)
\item
  time.Now() in bridge/perpetuals/teleport/attestation Run paths
  (consensus split) Re-review needed: yes -\/- C-01, C-02, C-03, H-01,
  H-02 Recommendation: do-not-ship
\end{enumerate}

\end{document}
